\section{Lecciones del método agéntico}
\label{sec:conclusiones:lecciones}

El Anexo~\ref{anx:uso-ia} declara qué tareas hicieron los agentes de inteligencia artificial y cuáles el autor.

Este trabajo separa a quien construye una Datasheet de quien la verifica.
El riesgo propio de un agente es dar por publicado un dato que la fuente vigente no contiene.
Le ocurrió a Claude Opus 5.5 en la prueba de la Sección~\ref{sec:comparacion:harvey}, que le planteó preguntas de comparación con las Datasheets en el prompt y sin ellas.
Sin ellas, atribuyó a Amazon SES en producción 50.000 correos al día y 14 por segundo.
La documentación vigente no da ninguna de las dos: la primera solo sale en una entrada del blog de AWS de 2014, y la segunda no está en ninguna fuente de su Datasheet.
El mismo error en un agente que construye una Datasheet pasaría al catálogo.

En este trabajo lo acotan dos controles.
El primero es que otro agente recibe la Datasheet terminada y la revisa campo a campo con las fuentes a la vista, en la fase de revisión del ciclo de la Sección~\ref{sec:propuesta:ciclo}.
El segundo es el verificador de citas, una comprobación automática que contrasta la cita literal de cada valor con la copia archivada de su página, y no se aplica a Postmark, cuyas páginas no se archivaron.
No se ha contado cuántos valores de las ocho Datasheets corrigieron esos dos controles, de modo que su eficacia no está medida.

El libro de criterios es el fichero de reglas común a las ocho Datasheets.
Cada duda de modelado que no cubría se decidió una sola vez y entró en él como regla nueva, de modo que no se volvió a plantear en otra API.
Así, los límites de recursos, como el número de dominios o de \textit{webhooks} de una cuenta, no se modelan como cuotas en ninguna de las ocho.

En la prueba de la Sección~\ref{sec:comparacion:harvey}, con las Datasheets en el prompt, Claude Opus 5.5 leyó bien los campos que pedían las cuatro preguntas, que tocan cinco de las ocho APIs.
Citó el campo de cada cifra, separó el valor publicado del interpretado y llegó a la misma conclusión que la referencia calculada a mano, la misma a la que llegó sin ellas.
La prueba comprueba la lectura y no que las Datasheets sean correctas: como se construyeron con la misma familia de modelos, un valor mal leído al construirlas pasaría a la respuesta sin que la prueba lo detectase.
Por eso un valor que solo ha deducido un agente se queda en ámbar, el color de lo interpretado, mientras no lo respalde un valor publicado o medido.

La reducción de horas del autor no se puede atribuir a los agentes.
Las horas suman \dato{total.real} frente a las \dato{total.prev} previstas, pero el tiempo de los agentes no se registró y no hay una medición del mismo trabajo sin ellos (Sección~\ref{sec:metodologia:desviaciones}).
Medir ese efecto exigiría registrar ese tiempo desde el inicio y construir sin agentes una muestra de Datasheets.
